\section{A complete 713-bit equational alternative}
\label{sec:e713}

This section gives a separate alternative to the binary machine in the main
report.  The supplied finite equational presentation has eight equations,
twelve operation symbols and two specified ground endpoints. Its entire
serialized consistency problem occupies 713 bits in the previously used
packet codec. Subject to the inherited ordinary calibration of raw EPSC,
we prove in the ZFC metatheory that
\begin{equation}
 \mathcal E_{713}\not\vdash N(d,d)=d
 \quad\longleftrightarrow\quad \operatorname{Con}(\mathrm{ZFC}).
 \label{eq:e713-calibration}
\end{equation}
Here derivability is ordinary equational derivability, including arbitrary
ground substitutions and contextual use in either direction. This is an
ordinary mathematical proof formalizable in ZFC, not a delivered Lean/Rocq
or internal PA proof. No binary Turing-machine state count is inferred from
713 bits. The 238-state construction in the main report uses the existing
calibrated machine and is independent of this alternative.

The audited numerical baseline is the complete 946-bit packet from the
19 October 2026 coordinated-reuse research archive. Our builder reconstructs its
equations, ordered signature, endpoints and bit string exactly. Source 79
supplies the earlier 1,169-bit native-stack presentation and its ordinary
model argument. The earlier 946-bit construction supplies the sharing
$J=P$, truth fusion and the allocation $\mathsf{Pick}(x)=W(x,d)$.
Boolean factoring and elimination of an auxiliary branch predicate are
standard transformations; we claim the explicit charged construction and
the new comparison with the inspected packet, without asserting exhaustive
historical priority. An earlier archive named
\texttt{zfc-sub240-research-20261019.zip} was described in conversation
history as containing a 766-bit kernel. Its bytes were unavailable in the
inspected material, so no literal comparison with that kernel is claimed.

\subsection{Raw EPSC and the exact inherited dependency}

Work with a nonempty first-order domain and a binary relation $E$. Define
\begin{align*}
 \mathsf{Empty}(x)&\;\Longleftrightarrow\;\forall u\,\neg E(u,x),\\
 y\subseteq_E z&\;\Longleftrightarrow\;\forall u\,(E(u,y)\to E(u,z)),\\
 \mathsf{LB}(R)&\;\Longleftrightarrow\;
       \forall x\,\exists z\,\forall y\,(R(x,y)\to y\subseteq_E z).
\end{align*}
For each equality-free pure-membership formula $R(x,y;\bar p)$ with any
finite parameter tuple, raw EPSC asserts that $\mathsf{LB}(R)$ implies
the existence of $h$ with
\begin{align}
 ((\mathsf{Empty}(x)\lor E(x,h))\land R(x,y))&\longrightarrow E(y,h),
 \label{eq:e713-close}\\
 E(y,h)&\longrightarrow\exists x\,R(x,y).
 \label{eq:e713-support}
\end{align}
These are actual formula instances, not quantification over every semantic
binary relation. The previously reviewed Sources 14 and 08 give the
ordinary ZFC-provable consistency comparison between this full raw schema
and ZFC. Source 14 first derives empty objects, power objects and profile
singletons, reconstructs transitive hulls and Separation/Union, and then
reduces EPSC to the earlier PSC schema. Source 08 supplies the hereditary
saturation interpretation and the ordinary Foundation/Choice comparison.
This section uses that calibration as an explicit inherited theorem.
It neither shortens those proofs nor assumes that an attractive incomplete
code calculus is already a complete proof system for the schema.

\subsection{Signature, native stack and the complemented binder}

The signature is
\[
 d^0,N^2,V^2,P^2,A^2,Q^1,W^2,
 \mathsf{Target}^1,\mathsf{Next}^2,\mathsf{Mem}^2,
 p_E^2,p_\varepsilon^2.
\]
Superscripts give arities. The first six equations are
\begin{align}
 N(N(x,N(N(y,x),x)),N(y,N(z,x)))&=y,\label{eq:e713-sh1}\\
 d&=N(d,N(d,d)),\label{eq:e713-top}\\
 V(P(x,s),d)&=x,\label{eq:e713-head}\\
 V(P(x,s),A(j,d))&=V(s,j),\label{eq:e713-tail}\\
 p_\varepsilon(s,P(a,b))&=N(p_\varepsilon(s,a),p_\varepsilon(s,b)),
 \label{eq:e713-nand}\\
 p_\varepsilon(s,A(i,j))&=p_E(V(s,i),V(s,j)).\label{eq:e713-atom}
\end{align}
Equation~\eqref{eq:e713-sh1} is the published Sheffer single axiom Sh1.
It makes the reduct a Boolean algebra with $N$ interpreted as NAND;
\eqref{eq:e713-top} makes $d$ its top. The native stack equations have
ordinary equality and apply to the entire domain. They are not identities
only after applying a truth-value map.

The seventh equation is the new binder:
\begin{equation}
\begin{split}
 p_\varepsilon(s,Q(a))={}&N\bigl(
      p_\varepsilon(P(W(s,a),s),a),\,
      p_\varepsilon(P(z,s),a)\bigr).
\end{split}
\label{eq:e713-quantifier}
\end{equation}
The variable $z$ is universally quantified, as are all equation variables.
The intended meaning of $Q$ is \emph{existential negation}:
\[
 \varepsilon(s,Q(a))\quad\longleftrightarrow\quad
      \exists z\,\neg\varepsilon(P(z,s),a).
\]
It is essential not to read it as the old ordinary existential constructor.

\paragraph{Complemented-witness lemma.}
Let $X$ be nonempty, let $(u_z)_{z\in X}$ be elements of a Boolean algebra,
and let $w=u_{z_0}$ be a selected member of that family. Then
\[
 (\forall z\;q=\neg(w\wedge u_z))
 \quad\longleftrightarrow\quad
 q=\neg w\ \land\ \forall z\;(w\leq u_z).
\]
Indeed, $z=z_0$ gives $q=\neg w$ by idempotence. Taking complements in
the other equations yields $w=w\wedge u_z$. The reverse implication
is immediate. In the two-element algebra this says exactly that $q$ is
true iff some $u_z$ is false, and the selected $u_{z_0}$ is then false.
The lemma does not say that every arbitrary Boolean-valued family has
such a selector: the selected value must be an attained minimum.

The new equation therefore replaces both the old selected-witness equality
and the arbitrary-witness implication. With the rest of the literal
946-bit packet retained, it gives 917 bits: one fewer equation and seven
fewer term nodes save $1+4\cdot7=29$ bits. This comparison does not depend
on the subsequent Boolean factoring or branch elimination.

\subsection{One fully expanded closure equation}

The following names abbreviate \emph{terms} for explanation only:
\begin{align*}
 h&=\mathsf{Target}(b),& n_z&=\mathsf{Next}(b,z),&m_z&=\mathsf{Mem}(b,z),\\
 R_b(x,y)&=p_\varepsilon(P(y,P(x,b)),V(b,d)),&&&
\end{align*}
and the eight Boolean-valued terms
\begin{align*}
 r&=R_b(h,n_z),&e&=p_E(m_z,n_z),&f&=p_E(m_z,z),\\
 a&=p_E(W(x,d),x),&b_0&=p_E(x,h),&c&=R_b(x,y),\\
 u&=p_E(y,h),&v&=R_b(\mathsf{Next}(b,y),y).&&
\end{align*}
The descriptor variable is $b$; $b_0$ is just the name of one predicate
value. Define $\overline t=N(d,t)$ and
\[
 \ell=N(d,N(c,N(a,N(d,b_0)))),\qquad
 B^-=N(r,N(d,N(e,N(d,f)))),
\]
\[
 H^-=N\bigl(N(\ell,N(d,u)),\,N(u,N(d,v))\bigr).
\]
The eighth equation is
\begin{equation}
                       d=N(B^-,H^-).\label{eq:e713-closure}
\end{equation}
Every name in this display is expanded in
\texttt{generated/E713.equations}, \texttt{generated/E713.json}, and
\texttt{generated/E713.bits}. No definition or arity is supplied to the
decoder from outside the packet.

In the Boolean algebra, write
\begin{align*}
 \mathsf{Bad}(b,z)&=r\land e\land\neg f,\\
 \mathsf{Hull}(b,x,y)&=(a\lor\neg c\lor u)
       \land(\neg b_0\lor\neg c\lor u)
       \land(\neg u\lor v).
\end{align*}
The elementary identities
\[
 \ell=c\land(\neg a\lor b_0),\qquad
 \mathsf{Hull}=(\ell\to u)\land(u\to v),
\]
give $B^-=\neg\mathsf{Bad}$ and $H^-=\neg\mathsf{Hull}$. Thus the
last equation is exactly
\begin{equation}
 \forall b,z,x,y\quad
       \mathsf{Bad}(b,z)\lor\mathsf{Hull}(b,x,y).\label{eq:e713-disj}
\end{equation}
All three former bad-branch literals and all three former hull clauses
remain present with their original polarities. Neither global support nor
the empty-input clause has been removed.

In two-valued logic, for each fixed $b$,~\eqref{eq:e713-disj} is equivalent
to the existence of a Boolean flag $g$ such that
\[
  (\forall z\;g\lor\mathsf{Bad}(b,z))
  \quad\land\quad
  (\forall x,y\;\neg g\lor\mathsf{Hull}(b,x,y)).
\]
If some $z$ makes the bad branch false, all hull instances must be true;
otherwise choose the bad branch. This explains the exact removal of
$p_G$ from the signature. This step is applied \emph{after} reducing
predicate values to two truth values. We do not assume that a general
Boolean algebra has the infinitary meet needed to manufacture such a
flag from all its predicate values.

\subsection{Lower interpretation and all binder translations}

Let a model of $\mathcal E_{713}$ separate $N(d,d)$ and $d$. Its
Sh1 reduct is a nontrivial Boolean algebra. Choose an ultrafilter and
apply its Boolean homomorphism to $\{0,1\}$ only to the values of
$p_E$ and $p_\varepsilon$. Keep the original domain and every native
function $P,V,A,Q,W,\mathsf{Target},\mathsf{Next},\mathsf{Mem}$.
The native equations retain their literal equality, and all Boolean
equations retain their logical consequences.

Equation~\eqref{eq:e713-quantifier} and the complemented-witness lemma
now give the stated $\exists\neg$ semantics. To compare with Source 79's
ordinary existential presentation, the complete operation translations are
\begin{align*}
 J_{\rm old}(a,b)&=P(a,b),&\mathsf{top}_{\rm old}&=d,\\
 Q_{\rm old}(a)&=Q(P(a,a)),&
 W_{\rm old}(s,a)&=W(s,P(a,a)),\\
 \mathsf{Pick}_{\rm old}(x)&=W(x,d).&&
\end{align*}
All other native functions are unchanged. For any fixed $s,a$, the class
$P(a,a)$ is the complement of $a$. Applying the new binder to it gives
\begin{align*}
 \varepsilon(s,Q(P(a,a)))&\longleftrightarrow
       \exists z\,\varepsilon(P(z,s),a),\\
 \varepsilon(s,Q(P(a,a)))&\longleftrightarrow
       \varepsilon(P(W(s,P(a,a)),s),a).
\end{align*}
The arbitrary-witness implication follows from the first line.
The independent old member selector is $W(x,d)$, not
$W(x,P(d,d))$. This is why we compare with Source 79 before its
member-selector fusion, rather than asserting an unjustified literal
expansion to the already-fused 946-bit equations.

For completeness, define selectors $\nu_0=d$ and
$\nu_{i+1}=A(\nu_i,d)$. Equations~\eqref{eq:e713-head}--\eqref{eq:e713-tail}
recover coordinate $i$ of every finite stack by $i$ tail substitutions
and one head substitution. An atomic formula on coordinates $i,j$ is
represented by $A(\nu_i,\nu_j)$, NAND by $P$, and ordinary existential
quantification of a represented matrix $a$ by $Q(P(a,a))$. Prepending
a stream entry shifts each free coordinate by one under a binder.
The descriptor has its code in its own coordinate 0, hence the full
environment $P(y,P(x,b))$ has output $y$ at 0, input $x$ at 1, code at 2,
and its finite parameter list starting at 3. Build the class term first
and then its enclosing descriptor. Structural induction represents every
finite equality-free pure-$E$ formula with every finite parameter tuple;
no satisfaction predicate for the expanded language is asserted.

Fix such a descriptor and suppose $\mathsf{LB}(R_b)$. Apply boundedness
at the object $h=\mathsf{Target}(b)$ to obtain a $z$ bounding every
$R_b(h,y)$-output. The three bad literals cannot all hold at this $z$:
they would give $R_b(h,n_z)$, $E(m_z,n_z)$ and $\neg E(m_z,z)$, contrary
to that bound. Equation~\eqref{eq:e713-disj} consequently gives every
hull clause. If $x$ is empty, $E(W(x,d),x)$ is false, so its first
clause gives closure at empty inputs. The second gives closure at members
of $h$. The last gives unrestricted support, witnessed by
$\mathsf{Next}(b,y)$. Thus every instance of raw EPSC holds.

The inherited raw-EPSC calibration now yields
$\operatorname{Con}(\mathrm{ZFC})$. No choice of an infinitary Boolean
flag and no quotient of native objects was needed.

\subsection{Upper model, including the free witness slice}

Assume $\operatorname{Con}(\mathrm{ZFC})$ while working in ZFC.
Completeness and downward L\"owenheim--Skolem give a countably infinite
set-coded ZFC model; relabel its domain by $\mathbb N$ and retain its
membership relation $E$. The model may be nonstandard and need not be
transitive. It supplies the required pure-membership EPSC instances by
the inherited upper calibration.

Take a Cantor pairing bijection $\pi$ with
$\pi(i,j)\geq\max(i,j)$ and set
\[
 d=0,\quad P(x,s)=3\pi(x,s)+1,\quad A(i,j)=3\pi(i,j)+2,
 \quad Q(a)=3a+3.
\]
The positive constructor ranges are disjoint and every constructor code
strictly exceeds its arguments. Interpret a $P$-code as a stream cell,
and give every non-$P$ stream head and tail 0. The index of a selector
is the number of successive outer layers $A(j,0)$; it is zero when the
selector has any other shape. Let $V(s,i)$ return that stream coordinate.
Both native stack equations hold for arbitrary, including noncanonical,
arguments.

Define class incidence by recursion on the natural-number class code:
$d$ is false; $P(a,b)$ is NAND of its child classes; $A(i,j)$ tests
$E(V(s,i),V(s,j))$; and $Q(a)$ means
$\exists z\,\neg\varepsilon(P(z,s),a)$. The strict decrease of child
codes makes this definition well founded. Every fixed class describes a
finite pure-$E$ formula in finitely many stream coordinates. Therefore
every fixed descriptor $b$ describes a pure-$E$ formula in $x,y$ with
finitely many parameters. The global external decoder is not claimed to
be internally definable in the chosen nonstandard set model.

For $a\ne d$, choose $W(s,a)$ to be the least external-index element
falsifying $\varepsilon(P(z,s),a)$ whenever one exists, and 0 otherwise.
At $a=d$, every instance is false and the new $Q(d)$ class is true.
Both sides of~\eqref{eq:e713-quantifier} then equal true, independently
of $W(s,d)$. Allocate this entire free slice as
\[
 W(x,d)=\begin{cases}
   \min\{u:E(u,x)\},&\text{if $x$ has an $E$-member},\\
   0,&\text{otherwise}.
 \end{cases}
\]
Thus the complemented binder preserves exactly the freedom needed for
the earlier member-selector allocation. No second, incompatible role is
assigned to that slice.

For each descriptor, if local boundedness fails choose $h$ with
unbounded output subsets. For each $z$ choose $n_z,m_z$ satisfying
$R_b(h,n_z)\land E(m_z,n_z)\land\neg E(m_z,z)$. Interpret
$\mathsf{Target},\mathsf{Next},\mathsf{Mem}$ by these choices. The bad
branch is true at every $z$, so the final equation holds. If local
boundedness holds, choose an EPSC hull $h$. For each $y\in_E h$, select
an input witnessing $R_b(\mathsf{Next}(b,y),y)$; use 0 elsewhere and
for the unused member witnesses. The member-selector slice makes the
first hull clause vacuous at nonempty $x$, and EPSC gives it at empty
$x$. The remaining two hull clauses follow directly from EPSC. Hence
the hull branch is always true. The two uses of $\mathsf{Next}$ are
compatible because they are assigned in mutually exclusive descriptor
cases. External least-index choices make all functions total.

Finally transport a countably infinite Boolean algebra, for example the
finite/cofinite algebra, to this same carrier with top equal to the
existing $d$. Interpret $N$ as its NAND and $p_E,p_\varepsilon$ as its
top or bottom according to the just-constructed relations. Keep every
native function fixed. All eight equations hold, and $N(d,d)\ne d$.
This proves the upper implication without confusing a Boolean quotient
with the native domain.

\subsection{Equational proof search and exact packet cost}

The ground terms exist because the signature contains $d$. Let $\sim$
be the congruence generated on all ground terms by every ground instance
of every equation. If the endpoints are not equationally derivable,
the ground-term quotient by $\sim$ is a separating model: finite choices
of representatives verify every universally quantified equation. The
converse is equational soundness. This ordinary set construction in ZFC
converts the model comparison above into~\eqref{eq:e713-calibration}.

Consequently a fair enumerator of finite ground contextual rewrite
certificates, stopping on a trace from $N(d,d)$ to $d$, has nonhalting
equivalent in ZFC to $\operatorname{Con}(\mathrm{ZFC})$. Reversed equations
may introduce variables absent from their matched side: each certificate
must specify ground values for the union of variables on both sides,
with shared values for repeated occurrences. This is a genuine
computable search specification, but its parser, substitution engine,
context traversal, enumeration, reset and blank-tape writer have not
been compiled into a charged binary machine here.

The inherited codec contains an 8-bit symbol count, three arity bits for
each symbol, an 8-bit variable capacity, one leading bit per equation,
a zero equation-list terminator and two ground endpoint trees. Symbols
and local variables share a fixed-width opcode space. For $\mathcal E_{713}$,
twelve symbols plus four variables require exactly four bits per node.
The complete count is
\[
 \underbrace{8+3\cdot12+8}_{52\text{ header}}
 +\underbrace{8+1+4\cdot159}_{645\text{ equations}}
 +\underbrace{4\cdot4}_{16\text{ endpoints}}=713.
\]
The literal eight equation-node counts are
$14,6,6,10,12,12,17,82$, totaling 159.
The bit-character SHA-256, excluding the file's final newline, is
\begin{center}\small\ttfamily
3cd0c3d0fb8d8640cfef1f44e0c22d377dd930de88a94093a2ef7650e22c9e72.
\end{center}

\begin{center}
\begin{tabular}{lrrrr}
\hline
Variant&Equations&Nodes&Symbols/variables&Complete bits\\\hline
Inspected shared-stack baseline&14&215&13/3&946\\
Complemented binder only&13&208&13/3&917\\
Factored closure with a flag&9&166&13/3&745\\
Factored closure without a flag&8&159&12/4&713\\\hline
\end{tabular}
\end{center}

The complete reduction relative to the inspected 946-bit packet is 233
bits, or about 24.63 percent. Every row includes its signature, all
equations and both endpoints; none is merely the size of an axiom basis.
The builder checks literal round trips, every strict-prefix truncation
and malformed/trailing inputs. It also checks the Boolean identities and
the attained-minimum binder lemma on explicit finite Boolean algebras.
These are finite transcription and algebra regressions, not models of
EPSC or tests of consistency. The accompanying independent review adds
its own parser and direct equation checks. The final receipt records the
exact executed scope.

\subsection{Further questions specific to this frontend}

The immediate engineering question is whether a ground-rewrite checker
specialized to these eight equations can be compiled more compactly than
the existing arithmetic proof-search engine. A useful implementation
must charge the entire 713-bit writer and the interfaces listed above.
The complemented-witness identity may also be applied to other class or
formula evaluators, but each adaptation must preserve its witness-slice
requirements and every occurrence of the old quantifier. Finally, an
explicit primitive-recursive proof translator for raw EPSC and the
complete equational kernel would strengthen the internal-arithmetic
calibration beyond the ordinary ZFC proof supplied here. These questions
remain separate from the achieved 238-state machine reduction.
